\documentclass[10pt]{article}
\usepackage[a4paper,margin=0.82in]{geometry}
\usepackage{amsmath,amssymb,bm}
\usepackage{booktabs}
\usepackage{xcolor}
\usepackage[colorlinks=true,urlcolor=blue]{hyperref}
\setlength{\parindent}{0pt}
\setlength{\parskip}{3pt}
\renewcommand{\arraystretch}{1.1}
\newcommand{\R}{\mathbf{R}}
\newcommand{\vd}{\mathbf{d}}
\newcommand{\vC}{\mathbf{C}}

\begin{document}
\begin{center}
{\large\bfseries Single--image georeferencing of a MicaSense frame over flat ground}\\[2pt]
{\small Mapping image pixels to ground coordinates $(E,N)$ from one capture --- no second image, no panel, no GCP}\\[3pt]
{\small\itshape Developed by Ali Lotfi, in collaboration with Claude Opus 4.8}
\end{center}
\vspace{-4pt}

\section*{1.\quad The idea}
An orthomosaic carries a 6--parameter \emph{affine} geotransform because orthorectification has already resampled the imagery onto a planar map grid, so pixel\,$\to$\,ground is linear. A \emph{raw} frame is a central (perspective) projection of the 3--D scene: a pixel defines a \emph{ray}, not a point, and a ground coordinate is obtained only by intersecting that ray with a surface. For an agricultural field the surface is well approximated by a plane, and then the pixel\,$\leftrightarrow$\,ground map of one frame is an \emph{exact} $3\times3$ projective homography (8 parameters) --- the projective generalisation of the ortho's affine. A tilt away from nadir contributes the off--diagonal (keystone) terms that an affine cannot represent.

\section*{2.\quad Camera model from the image metadata}
Every parameter below is read directly from the file's XMP/EXIF.

\noindent\textbf{Intrinsics (pixels).} With $\rho=$ \texttt{FocalPlaneXResolution} (px\,mm$^{-1}$), focal length $f=f_{\mathrm{mm}}\rho$ and principal point $(c_x,c_y)=(\mathrm{pp}_x\rho,\,\mathrm{pp}_y\rho)$, giving
\[
\mathbf{K}=\begin{bmatrix} f & 0 & c_x\\ 0 & f & c_y\\ 0 & 0 & 1 \end{bmatrix}.
\]
We use the OpenCV camera frame: $x$ right, $y$ down, $z$ along the optical axis (downward at nadir).

\noindent\textbf{Lens distortion.} Brown model, radial $k_1,k_2,k_3$ and tangential $p_1,p_2$ (\texttt{Camera:PerspectiveDistortion}, in that order).

\noindent\textbf{Exterior orientation.} The projection centre is $\vC=(E_0,N_0,h)^{\!\top}$, where $(E_0,N_0)$ is the GPS longitude/latitude projected into the local UTM zone and $h$ the GPS altitude. Note: MicaSense records \texttt{GPSAltitude} as height above the WGS84 \emph{ellipsoid}; this is automatically converted to orthometric (MSL) height via the EGM96 geoid grid so that it is comparable with DEM ground elevations. The orientation is the rotation $\R$ that takes camera axes into the world East--North--Up (ENU) frame, built from the EXIF Yaw/Pitch/Roll $(\psi,\theta,\varphi)$:
\[
\R \;=\; \mathbf{C}_{en}\,\mathbf{C}_{nb}(\psi,\theta,\varphi)\,\mathbf{R}_{cb},
\qquad
\mathbf{R}_{cb}=\!\begin{bmatrix}0&-1&0\\1&0&0\\0&0&1\end{bmatrix}\!,\quad
\mathbf{C}_{en}=\!\begin{bmatrix}0&1&0\\1&0&0\\0&0&-1\end{bmatrix}\!.
\]
Here $\mathbf{R}_{cb}$ is the fixed look--down camera$\to$body map, $\mathbf{C}_{en}$ swaps the navigation NED frame to ENU, and $\mathbf{C}_{nb}=\mathbf{R}_z(\psi)\mathbf{R}_y(\theta)\mathbf{R}_x(\varphi)$ is the body$\to$NED rotation (B\"aumker convention, $\psi$ yaw, $\theta$ pitch, $\varphi$ roll).

\section*{3.\quad From a pixel to ground coordinates $(E,N)$}
For a pixel $(u,v)$:

\textbf{(i) Normalise:} $\;x_d=(u-c_x)/f,\quad y_d=(v-c_y)/f.$

\textbf{(ii) Undistort} --- recover the ideal $(x,y)$ by inverting the Brown model, with $r^2=x^2+y^2$,
\[
x_d = x\bigl(1+k_1r^2+k_2r^4+k_3r^6\bigr)+2p_1xy+p_2(r^2+2x^2),
\]
\[
y_d = y\bigl(1+k_1r^2+k_2r^4+k_3r^6\bigr)+p_1(r^2+2y^2)+2p_2xy,
\]
solved by a few fixed--point iterations (MicaSense distortion is sub--pixel except at the edges).

\textbf{(iii) Form the world ray:} the camera--frame direction $\vd_c=(x,y,1)^{\!\top}$ rotates to $\vd=\R\,\vd_c$.

\textbf{(iv) Intersect the ground plane} $Z=Z_g$ (with $Z_g$ the field elevation). The ray is $\mathbf{P}(t)=\vC+t\,\vd$; the plane fixes the scale,
\[
\boxed{\;t^{\*}=\dfrac{Z_g-C_z}{d_z},\qquad
\begin{bmatrix}E\\N\end{bmatrix}=
\begin{bmatrix}C_x\\C_y\end{bmatrix}+t^{\*}\begin{bmatrix}d_x\\d_y\end{bmatrix}.\;}
\]
That pair $(E,N)$ is the ground coordinate of the pixel, in UTM.

\textbf{Equivalent closed form (the matrix).} For a plane the entire map is a single homography. Place the local ground origin at the nadir point (plane $Z=0$, camera at $(0,0,H_{\!A})$ with $H_{\!A}=h-Z_g$). Then
\[
s\,(u,v,1)^{\!\top}=\mathbf{K}\,[\,\mathbf{r}_1\ \ \mathbf{r}_2\ \ \mathbf{t}\,]\,(E,N,1)^{\!\top},
\qquad \mathbf{t}=-H_{\!A}\,\mathbf{r}_3,
\]
where $\mathbf{r}_1,\mathbf{r}_2,\mathbf{r}_3$ are the columns of $\R^{\!\top}$ (world$\to$camera). Inverting the $3\times3$ matrix gives $(E,N)$ directly from $(u,v)$; absolute coordinates are this plus $(E_0,N_0)$. The relation is exact for a plane and applies to \emph{undistorted} pixels (do step ii first).

\section*{4.\quad What to expect}
Because the model is geometrically exact for a planar surface, pixel\,$\leftrightarrow$\,ground round--trips to numerical zero. Three practical limits remain. \textbf{(a) Vertical datum:} the GPSAltitude in the metadata is an ellipsoid height (WGS84); DEM ground elevations are orthometric (above the geoid/MSL). The script converts automatically using EGM96 so the two are in the same datum. If the geoid grid is unavailable, the fallback is to treat GPS as orthometric with a warning. For absolute-certainty workflows, give the AGL directly via \verb|--agl| to sidestep the datum question entirely. \textbf{(b) Absolute position} inherits the onboard GNSS uncertainty --- roughly $1$--$2$\,m for a standard fix (the EXIF \texttt{GPSXYAccuracy}/\texttt{GPSZAccuracy} report it), removable only with RTK/PPK or a ground control point. \textbf{(c) Canopy parallax:} at low flight height a differential crop canopy behaves like relief and shifts edge pixels by about $(h_{\mathrm{canopy}}/H_{\!A})\times(\text{edge distance})$ --- decimetre--scale for a $1$\,m canopy at $22$\,m AGL, vanishing toward nadir. A \emph{uniform} field slope is harmless, since a homography handles any tilted plane.

\section*{5.\quad Running it}
Dependencies: \texttt{numpy}, \texttt{pyproj}, \texttt{Pillow}, \texttt{tifffile} (and \texttt{rasterio} only for GeoTIFF output).

\noindent Print the homography and self--checks:
\begin{quote}\small\verb|python micasense_georef.py IMG_0010_4.tif --ground-elev 481.5|\end{quote}

\noindent Write a north--up GeoTIFF (UTM zone auto--selected from the GPS):
\begin{quote}\small\verb|python micasense_georef.py IMG_0010_4.tif --ground-elev 481.5 --gsd 0.012 --out out.tif|\end{quote}

\noindent Alternatively, give the camera height above ground directly (sidesteps vertical-datum ambiguity):
\begin{quote}\small\verb|python micasense_georef.py IMG_0010_4.tif --agl 22.0 --out out.tif|\end{quote}

\noindent\verb|--ground-elev| is the field elevation in metres (orthometric/MSL from a DEM; it sets the footprint's scale and position), and \verb|--gsd| the output pixel size in metres. \verb|--agl| is an alternative that gives the absolute height above ground in metres (from the flight plan or other source); exactly one of \verb|--ground-elev| or \verb|--agl| is required. Run the script per band \emph{after} band--to--band alignment to build a georeferenced multispectral stack.
\end{document}
